\documentclass[11pt,a4paper]{article}
\usepackage[margin=1in]{geometry}
\usepackage{amsmath,amssymb,amsthm}
\usepackage{algorithm}
\usepackage{algpseudocode}
\usepackage{booktabs}
\usepackage{hyperref}

\theoremstyle{definition}
\newtheorem{definition}{\textbf{Definition}}[section]
\newtheorem{theorem}{\textbf{Theorem}}[section]
\newtheorem{lemma}[theorem]{\textbf{Lemma}}
\newtheorem{remark}{\textbf{Remark}}[section]

\title{\textbf{Direct Preference Optimization (DPO)}}
\author{\textbf{Team Scaibu} \\ \small Email: \texttt{jaydeep.9664920749@gmail.com} \\ \small \textbf{Architecture and Core Engine Team}}
\date{\textbf{October 2026}}

\begin{document}
\maketitle

\section{Definitions and Cognitive Mental Models}

\subsection{Formal Mathematical Definition}
\textbf{Definition (Implicit Reward Reparameterization and Direct Policy Alignment):}
Traditional Reinforcement Learning from Human Feedback (RLHF) optimizes policies via a two-stage process: first fitting an explicit reward model $r_\psi(x, y)$ to pairwise preference data via maximum likelihood, and second optimizing the language model policy $\pi_\theta$ against $r_\psi$ via complex actor-critic reinforcement learning (PPO). \textbf{Direct Preference Optimization (DPO)} establishes that the constrained RLHF objective can be solved exactly in closed form, eliminating the reward model, the actor-critic loops, and the dynamic sampling entirely.

Under the standard KL-constrained RLHF objective:
{\boldmath
\begin{equation}
\max_\pi \mathbb{E}_{x \sim \mathcal{D}, y \sim \pi}[r(x, y)] - \beta D_{\text{KL}}(\pi(y \mid x) \,||\, \pi_{\text{ref}}(y \mid x))
\end{equation}
}
the analytically optimal policy satisfies $\pi^*(y \mid x) = \frac{1}{Z(x)} \pi_{\text{ref}}(y \mid x) \exp\left(\frac{1}{\beta} r(x, y)\right)$. Inverting this relationship reveals that the ground-truth reward can be expressed identically through the policy's log-ratio:
{\boldmath
\begin{equation}
r(x, y) = \beta \ln \frac{\pi^*(y \mid x)}{\pi_{\text{ref}}(y \mid x)} + \beta \ln Z(x)
\end{equation}
}
Substituting this \textbf{Implicit Reward Formulation} into the Bradley-Terry preference likelihood $P(y_w \succ y_l \mid x) = \sigma(r(x, y_w) - r(x, y_l))$ causes the intractable partition function $Z(x)$ to cancel out completely. The resulting \textbf{DPO Objective} optimizes the policy parameters $\theta$ directly over static pairwise comparison data:
{\boldmath
\begin{equation}
\mathcal{L}_{\text{DPO}}(\theta; \pi_{\text{ref}}) = - \mathbb{E}_{(x, y_w, y_l) \sim \mathcal{D}} \left[ \ln \sigma\left( \beta \ln \frac{\pi_\theta(y_w \mid x)}{\pi_{\text{ref}}(y_w \mid x)} - \beta \ln \frac{\pi_\theta(y_l \mid x)}{\pi_{\text{ref}}(y_l \mid x)} \right) \right]
\end{equation}
}

\subsection{Expanded Non-Technical Definition (Intuition for Anyone)}
\textbf{Intuitive Conceptual Definition:}
\begin{enumerate}
    \item \textbf{Cutting Out the Middleman}: In traditional RLHF, you build a robot judge (the reward model), and then use RL to teach the student model how to impress that judge. DPO realizes the judge is completely unnecessary! The math allows the student model to learn directly from the human preference votes in a single simple step.
    \item \textbf{Implicit Rewards}: Although DPO has no separate reward model, it secretly computes an internal score: ``How much more likely did I make this sentence compared to my old reference self?'' ($r_{\text{implicit}} = \beta \ln(\pi / \pi_{\text{ref}})$).
    \item \textbf{Zero Hallucination of Rewards}: Because DPO trains on a fixed dataset of pairs with standard supervised-style cross-entropy, it never suffers from PPO reward hacking, training collapse, or critic divergence.
    \item \textbf{Massive Hardware Savings}: Instead of loading 4 separate neural networks onto multiple GPU servers, DPO only needs the active model and a frozen reference, running over $3\times$ faster on standard hardware.
\end{enumerate}

\subsection{Child-Friendly Mental Model (Physical Metaphor)}
Imagine learning to bake cakes:
\begin{itemize}
    \item \textbf{The Old Method (PPO)}: You hire a food critic, teach the food critic what cake tastes like, and then bake 10,000 cakes in your kitchen while the critic grades each one.
    \item \textbf{The DPO Shortcut}: Your grandmother hands you two cakes side-by-side: ``The chocolate cake on the left is delicious; the burnt cake on the right is bad.'' You simply look at the recipes, increase the chocolate and reduce the burn!
    \item \textbf{Result}: You master the winning recipe in one afternoon without ever needing to hire a food critic!
\end{itemize}

\section{Core Mathematical Equations and Definitions}

\subsection{Implicit Reward Definition}
{\boldmath
\begin{equation}
\hat{r}_\theta(x, y) = \beta \ln \frac{\pi_\theta(y \mid x)}{\pi_{\text{ref}}(y \mid x)}
\end{equation}
}
\begin{itemize}
    \item \textbf{Formal Definition}: The closed-form implicit reward function mapping policy likelihood ratios to scalar utility.
    \item \textbf{What It Means}: High scores are assigned to responses that the active policy considers substantially more probable than the base reference model.
    \item \textbf{Why It Is Important}: Allows evaluating policy rewards and margins without needing a separate reward model architecture.
\end{itemize}

\subsection{DPO Pairwise Loss Function}
{\boldmath
\begin{equation}
\mathcal{L}_{\text{DPO}} = - \frac{1}{B}\sum_{i=1}^B \ln \sigma\left( \beta \left[ \left( \ln \pi_\theta(y_{w, i}) - \ln \pi_{\text{ref}}(y_{w, i}) \right) - \left( \ln \pi_\theta(y_{l, i}) - \ln \pi_{\text{ref}}(y_{l, i}) \right) \right] \right)
\end{equation}
}
\begin{itemize}
    \item \textbf{Formal Definition}: Binary classification loss maximizing the margin between chosen and rejected response likelihoods.
    \item \textbf{What It Means}: Increases the probability of the preferred response relative to the rejected response, anchored to the reference model.
    \item \textbf{Why It Is Important}: The most widely adopted preference optimization algorithm in modern open-weight LLMs (e.g. Zephyr, Llama-3-Instruct).
\end{itemize}

\subsection{DPO Policy Gradient Weighting}
{\boldmath
\begin{equation}
\nabla_\theta \mathcal{L}_{\text{DPO}} = - \beta \mathbb{E}\left[ \sigma(\hat{r}_\theta(x, y_l) - \hat{r}_\theta(x, y_w)) \left( \nabla_\theta \ln \pi_\theta(y_w \mid x) - \nabla_\theta \ln \pi_\theta(y_l \mid x) \right) \right]
\end{equation}
}
\begin{itemize}
    \item \textbf{Formal Definition}: Gradient update vector weighted by the misclassification probability of the current implicit reward margin.
    \item \textbf{What It Means}: Applies large gradient pushes to pairs where the model currently prefers the wrong response, and small updates to pairs already correctly ordered.
    \item \textbf{Why It Is Important}: Provides automatic dynamic sample weighting analogous to AdaBoost and hard-negative mining.
\end{itemize}

\section{Rigorous Mathematical Analysis Checklist}

\begin{itemize}
    \item \textbf{Notation}: $\pi_\theta$ (trained policy), $\pi_{\text{ref}}$ (frozen base model), $y_w$ (chosen completion), $y_l$ (rejected completion), $\beta$ (regularization scale).
    \item \textbf{Epistemic Status}: Discovered by Rafailov, Sharma, Mitchell, Ermon, Manning, and Finn (Stanford, NeurIPS 2023).
    \item \textbf{Dependency Graph}: Preference Data $(x, y_w, y_l) \to \text{Evaluate } (\ln \pi_\theta, \ln \pi_{\text{ref}}) \to \text{Implicit Margin} \to \nabla_\theta \mathcal{L}_{\text{DPO}}$.
    \item \textbf{Dimensions}: Batch size $B$, token sequence lengths $|y_w|, |y_l|$.
    \item \textbf{Regimes}: Instruction alignment, safety tuning, preference distillation, conversational formatting.
    \item \textbf{Invariants}: When $\pi_\theta \equiv \pi_{\text{ref}}$ on all candidate pairs, $\hat{r}_\theta \equiv 0.0$ and $\mathcal{L}_{\text{DPO}} \equiv \ln(2) \approx 0.693147$.
    \item \textbf{Isomorphisms}: DPO is mathematically isomorphic to weighted contrastive divergence on conditional energy-based models.
\end{itemize}

\section{Theoretical Theorems and Formal Proofs}

\begin{theorem}[Exact Derivation of the DPO Objective from the RLHF Dual]
The Direct Preference Optimization loss $\mathcal{L}_{\text{DPO}}$ is an exact mathematical equivalent to maximum likelihood training under the Bradley-Terry preference model parameterized by the optimal RLHF policy.
\end{theorem}

\begin{proof}
Consider the KL-regularized RLHF objective for a given prompt $x$:
{\boldmath
\begin{equation}
\max_\pi \sum_y \pi(y \mid x) r(x, y) - \beta \sum_y \pi(y \mid x) \ln \frac{\pi(y \mid x)}{\pi_{\text{ref}}(y \mid x)}
\end{equation}
}
Subject to the simplex constraint $\sum_y \pi(y \mid x) = 1$, we formulate the Lagrangian:
{\boldmath
\begin{equation}
\mathcal{L}(\pi, \lambda) = \sum_y \pi(y \mid x) \left[ r(x, y) - \beta \ln \frac{\pi(y \mid x)}{\pi_{\text{ref}}(y \mid x)} \right] - \lambda \left( \sum_y \pi(y \mid x) - 1 \right)
\end{equation}
}
Taking the functional derivative with respect to $\pi(y \mid x)$ and setting to zero:
{\boldmath
\begin{align}
\frac{\partial \mathcal{L}}{\partial \pi(y \mid x)} &= r(x, y) - \beta \left( \ln \frac{\pi(y \mid x)}{\pi_{\text{ref}}(y \mid x)} + 1 \right) - \lambda = 0 \\
\ln \frac{\pi(y \mid x)}{\pi_{\text{ref}}(y \mid x)} &= \frac{1}{\beta} r(x, y) - \left( \frac{\lambda}{\beta} + 1 \right) \\
\pi^*(y \mid x) &= \pi_{\text{ref}}(y \mid x) \exp\left( \frac{1}{\beta} r(x, y) \right) \exp\left( -\frac{\lambda}{\beta} - 1 \right)
\end{align}
}
Enforcing the normalization constraint defines the partition function $Z(x) = \sum_{y'} \pi_{\text{ref}}(y' \mid x) \exp\left(\frac{1}{\beta} r(x, y')\right) = \exp\left(\frac{\lambda}{\beta} + 1\right)$.
Taking logarithms on both sides:
{\boldmath
\begin{equation}
\ln \pi^*(y \mid x) = \ln \pi_{\text{ref}}(y \mid x) + \frac{1}{\beta} r(x, y) - \ln Z(x)
\end{equation}
}
Rearranging to isolate the reward function $r(x, y)$:
{\boldmath
\begin{equation}
r(x, y) = \beta \ln \frac{\pi^*(y \mid x)}{\pi_{\text{ref}}(y \mid x)} + \beta \ln Z(x)
\end{equation}
}
Now substitute this expression into the Bradley-Terry preference probability:
{\boldmath
\begin{align}
P(y_w \succ y_l \mid x) &= \sigma\left( r(x, y_w) - r(x, y_l) \right) \\
&= \sigma\left( \left( \beta \ln \frac{\pi^*(y_w \mid x)}{\pi_{\text{ref}}(y_w \mid x)} + \beta \ln Z(x) \right) - \left( \beta \ln \frac{\pi^*(y_l \mid x)}{\pi_{\text{ref}}(y_l \mid x)} + \beta \ln Z(x) \right) \right) \\
&= \sigma\left( \beta \ln \frac{\pi^*(y_w \mid x)}{\pi_{\text{ref}}(y_w \mid x)} - \beta \ln \frac{\pi^*(y_l \mid x)}{\pi_{\text{ref}}(y_l \mid x)} \right)
\end{align}
}
Notice that the uncomputable partition function $\beta \ln Z(x)$ cancels identically.
Minimizing the negative log-likelihood $-\ln P(y_w \succ y_l \mid x)$ over observed comparisons yields the exact DPO objective, completing the proof.
\end{proof}

\begin{theorem}[Vanishing Gradient Property on Saturated Margins]
Let the implicit reward margin be $h = \hat{r}_\theta(y_w) - \hat{r}_\theta(y_l)$. As the margin $h \to +\infty$, the gradient norm $\|\nabla_\theta \mathcal{L}_{\text{DPO}}\| \to 0$ exponentially.
\end{theorem}

\begin{proof}
Let $\mathcal{L}(h) = -\ln \sigma(h)$.
The derivative with respect to parameter $\theta$ is:
{\boldmath
\begin{equation}
\nabla_\theta \mathcal{L}(h) = -\frac{\sigma'(h)}{\sigma(h)} \nabla_\theta h = -(1 - \sigma(h)) \nabla_\theta h
\end{equation}
}
Recall that $1 - \sigma(h) = \sigma(-h) = \frac{1}{1 + e^h}$.
Taking the norm of the gradient:
{\boldmath
\begin{equation}
\|\nabla_\theta \mathcal{L}_{\text{DPO}}\| = \frac{1}{1 + e^h} \|\nabla_\theta h\| \le e^{-h} \|\nabla_\theta h\|
\end{equation}
}
As $h \to +\infty$, $e^{-h} \to 0$ exponentially fast.
Therefore, once the policy establishes a confident reward margin on a pair, gradient updates naturally vanish, preventing over-training on easy preference comparisons.
\end{proof}

\section{Foundational Architectural Questions and Epistemic Meta-Questions}

\subsection{Architectural Question 1: Likelihood Displacement Failure Mode}
\textbf{Question:} What is the ``likelihood displacement'' pathology in DPO, and how is it mitigated?\\
\textbf{Answer:} While DPO successfully increases the relative margin between chosen and rejected responses, it can achieve this either by increasing $P(y_w)$ or by driving $P(y_l)$ down to zero. In practice, models frequently lower the absolute probability of both responses while making $P(y_l)$ drop faster, degrading overall generation fluency. Adding a supervised fine-tuning loss regularizer $\mathcal{L}_{\text{DPO}} + \alpha \mathcal{L}_{\text{SFT}}(y_w)$ prevents likelihood displacement.

\textbf{Meta-Question:} How should the scaling parameter $\beta$ be chosen?\\
\textbf{Answer:} The parameter $\beta$ acts as an inverse temperature (typically $\beta \in [0.05, 0.2]$). Small $\beta$ ($0.01$) allows large divergence from the reference model, risking hallucinations and formatting degradation. Large $\beta$ ($0.5$) anchors the model rigidly to the base model, stifling preference alignment.

\subsection{Architectural Question 2: Online vs Offline DPO}
\textbf{Question:} Why is Online DPO (iterative DPO) increasingly replacing static DPO?\\
\textbf{Answer:} Static DPO trains on a fixed offline dataset. As the policy improves, it generates new styles and errors that were never present in the initial offline comparisons. Online DPO samples new completions from the active policy every $K$ iterations, annotates them with a verifier or judge model, and retrains, closing the distribution shift gap.

\textbf{Meta-Question:} Can DPO be applied without a frozen reference model?\\
\textbf{Answer:} Yes. Reference-free variants like SimPO and ORPO eliminate the reference model by substituting length-normalized average token log-likelihoods or odds ratios, saving memory while eliminating reference compute overhead.

\section{Algorithmic Implementation Specification}

\begin{algorithm}[H]
\caption{DPO Loss Evaluation and Implicit Reward Tracking}
\begin{algorithmic}[1]
\Require Policy log-probs $\mathbf{lp}_w, \mathbf{lp}_l$, reference log-probs $\mathbf{ref}_w, \mathbf{ref}_l$, scaling factor $\beta$
\Ensure DPO loss $\mathcal{L}_{\text{DPO}}$, implicit rewards $\mathbf{r}_w, \mathbf{r}_l$, mean margin $\bar{\Delta r}$, accuracy $\operatorname{Acc}$
\State Initialize $loss\_accum \gets 0.0, \quad correct \gets 0, \quad margin\_accum \gets 0.0$
\State Allocate $\mathbf{r}_w \in \mathbb{R}^B, \quad \mathbf{r}_l \in \mathbb{R}^B$
\For{$i = 0$ \textbf{to} $B-1$}
    \State $rw \gets \beta \cdot (\mathbf{lp}_w[i] - \mathbf{ref}_w[i])$
    \State $rl \gets \beta \cdot (\mathbf{lp}_l[i] - \mathbf{ref}_l[i])$
    \State $\mathbf{r}_w[i] \gets rw, \quad \mathbf{r}_l[i] \gets rl$
    \State $margin \gets rw - rl$
    \State $margin\_accum \gets margin\_accum + margin$
    \If{$margin > 0.0$}
        \State $correct \gets correct + 1$
    \EndIf
    \If{$margin \ge 0.0$}
        \State $loss\_accum \gets loss\_accum + \ln(1 + \exp(-margin))$
    \Else
        \State $loss\_accum \gets loss\_accum + (-margin + \ln(1 + \exp(margin)))$
    \EndIf
\EndFor
\State $\mathcal{L}_{\text{DPO}} \gets loss\_accum / B, \quad \bar{\Delta r} \gets margin\_accum / B, \quad \operatorname{Acc} \gets correct / B$
\State \Return $(\mathcal{L}_{\text{DPO}}, \mathbf{r}_w, \mathbf{r}_l, \bar{\Delta r}, \operatorname{Acc})$
\end{algorithmic}
\end{algorithm}

\section{Big-$\Theta$ Complexity Analysis}

\begin{itemize}
    \item \textbf{Implicit Reward Computation Time}: $\Theta(B)$ scalar multiplications and subtractions.
    \item \textbf{Softplus Loss Evaluation Time}: $\Theta(B)$ numerically stable softplus steps.
    \item \textbf{Preference Accuracy Tracking}: $\Theta(B)$ comparisons.
    \item \textbf{Space Complexity}: $\Theta(B)$ buffer for reward arrays.
    \item \textbf{Parallel Reduction Depth}: $\mathcal{D} = \Theta(\log B)$ for batch mean loss reduction.
\end{itemize}

\section{Single Canonical Repository Reference}

The complete deterministic scalar implementation, verification suite, and unit test benchmarks are published at:
\begin{center}
\url{https://github.com/Chief-Strategist-J/llm-obs-infra}
\end{center}

\end{document}
